\documentclass[12pt, a4paper]{article}
\usepackage[margin=1in]{geometry}
\usepackage{amsmath}
\usepackage{xcolor}
\usepackage{listings}
\usepackage{hyperref}

% C++ code formatting
\lstset{
    language=C++,
    basicstyle=\ttfamily\small,
    keywordstyle=\color{blue}\bfseries,
    stringstyle=\color{red},
    commentstyle=\color{teal},
    morecomment=[l][\color{magenta}]{\#},
    numbers=left,
    numberstyle=\tiny\color{gray},
    stepnumber=1,
    breaklines=true,
    frame=single,
    showstringspaces=false,
    tabsize=4
}

\title{\textbf{Lecture 1: Time Complexity \& C++ STL Basics}}
\author{Basics of CP Track}
\date{}

\begin{document}

\maketitle

\section{Counting Operations \& Time Complexity}

\subsection{Brute Force Example: Two Sum}
Given an array of $N$ integers, find if there are two numbers that add up to a target sum $X$.

\begin{lstlisting}
int count = 0;
for (int i = 0; i < n; i++) {
    for (int j = i + 1; j < n; j++) {
        if (arr[i] + arr[j] == X) {
            count++;
        }
    }
}
\end{lstlisting}

\textbf{How many operations does this take?}
\begin{itemize}
    \item When $i = 0$, inner loop runs $N-1$ times.
    \item When $i = 1$, inner loop runs $N-2$ times.
    \item Total = $(N-1) + (N-2) + \dots + 1 = \frac{N(N-1)}{2} \approx O(N^2)$.
\end{itemize}

\subsection{The CP Rule of Thumb}
\begin{center}
    \fbox{\textbf{A modern judge can perform roughly $10^8$ operations per second.}}
\end{center}
If $N = 10^5$, an $O(N^2)$ solution takes $\approx 10^{10}$ operations, requiring $\sim 100$ seconds. This will result in a \textbf{Time Limit Exceeded (TLE)}. We need algorithms closer to $O(N \log N)$ or $O(N)$. \textit{We will learn techniques ahead that help solve this problem much faster!}

\subsection{Linear vs. Binary Search}
Searching for a number in a \textbf{sorted} array of size $N = 10^6$:
\begin{itemize}
    \item \textbf{Linear Search:} Checks every element. Worst case: $10^6$ operations ($O(N)$).
    \item \textbf{Binary Search:} Halves the search space every step. Worst case: $\log_2(10^6) \approx 20$ operations ($O(\log N)$). 
\end{itemize}

\vspace{0.5cm}
\hrule
\vspace{0.5cm}

\section{The C++ STL (Standard Template Library)}

\subsection{\texttt{vector} and Iterators}
A dynamic array.
\begin{itemize}
    \item \texttt{push\_back(x)}: Adds to the end ($O(1)$).
    \item \texttt{v[i]}: Access element at index $i$ ($O(1)$).
    \item \texttt{pop\_back()}: Removes the last element ($O(1)$).
    \item \texttt{insert()}: Inserts at a position ($O(N)$ - \textit{Avoid inside loops!}).
\end{itemize}

\textbf{Iterators:} Pointers to elements in the container.
\begin{itemize}
    \item \texttt{v.begin()}: Points to the first element.
    \item \texttt{v.end()}: Points to \textit{one past} the last element.
\end{itemize}

\textbf{Sorting \& Binary Search (Requires sorted array):}
\begin{lstlisting}
sort(v.begin(), v.end()); // O(N log N)
auto it1 = lower_bound(v.begin(), v.end(), val); // First element >= val
auto it2 = upper_bound(v.begin(), v.end(), val); // First element > val
\end{lstlisting}


\subsection{\texttt{stack}}
LIFO (Last-In, First-Out) structure.
\begin{itemize}
    \item \texttt{push(x)}: Add to top ($O(1)$).
    \item \texttt{pop()}: Remove from top ($O(1)$).
    \item \texttt{top()}: Look at top element ($O(1)$).
\end{itemize}
\textit{Application Idea:} Finding the \textbf{Nearest Smaller Values} (Monotonic Stack) - a classic CSES problem.

\subsection{\texttt{queue}}
FIFO (First-In, First-Out) structure.
\begin{itemize}
    \item \texttt{push(x)}: Add to back ($O(1)$).
    \item \texttt{pop()}: Remove from front ($O(1)$).
    \item \texttt{front()}: Look at front element ($O(1)$).
\end{itemize}
\textit{Note: We won't do practice problems for this today. We will revisit queues extensively when we reach Breadth-First Search (BFS) in the Graphs section.}

\subsection{Pairs \& Custom Comparators}
Useful for binding two pieces of data, like coordinates or an element and its original index.

\begin{lstlisting}
vector<pair<int, int>> v;
v.push_back({5, 2});
// By default, sorts by first element, then second.
\end{lstlisting}

\textbf{Custom Comparator Example (Sorting pairs by the second element):}
\begin{lstlisting}
bool comp(const pair<int, int>& a, const pair<int, int>& b) {
    return a.second < b.second; 
}
sort(v.begin(), v.end(), comp);
\end{lstlisting}
\textbf{Crucial Rule:} Your comparator \textbf{must} return \texttt{false} if the elements are equal (i.e., \texttt{comp(a, a) == false}). Using \texttt{<=} will cause a Segmentation Fault!

\subsection{Set, Multiset, and Map}
Under the hood, these are balanced Binary Search Trees. Operations are $O(\log N)$.
\begin{itemize}
    \item \textbf{\texttt{set}:} Maintains a sorted collection of \textit{unique} elements.
    \item \textbf{\texttt{multiset}:} Maintains a sorted collection, allows \textit{duplicates}.
    \item \textbf{\texttt{map}:} Key-value pairs sorted by key. Useful for frequency counting (\texttt{freq[x]++}).
\end{itemize}

\vspace{0.5cm}
\hrule

\section{C++ Syntax Survival Guide}
These are the "obvious" syntax tricks that competitive programmers use everywhere, but can look like magic to beginners.

\subsection{Range-Based For Loops}
Instead of writing \texttt{for(int i = 0; i < v.size(); i++)}, C++11 introduced a much cleaner way to iterate through collections:

\begin{lstlisting}
vector<int> v = {1, 2, 3, 4, 5};

// 1. By Value (Makes a copy, safe but slower for big data)
for (int x : v) {
    cout << x << " "; 
}

// 2. By Reference (Avoids copies, allows modifying original data)
for (auto& x : v) {
    x *= 2; // This permanently changes elements in 'v'
}
\end{lstlisting}

\subsection{Iterators: Referencing and Incrementing}
An iterator is essentially a pointer to an element inside a container. We use the \texttt{auto} keyword to avoid typing massive types like \texttt{vector<int>::iterator}.

\begin{lstlisting}
auto it = v.begin(); // Points to the first element

// Dereferencing: Use the '*' operator to get the actual value
cout << *it << "\n"; 

// Incrementing: Move the iterator to the next element
it++; 
\end{lstlisting}

\textbf{Crucial Trap for Sets and Maps:} 
For a \texttt{vector}, you can jump around by doing \texttt{it += 5} (random access). For a \texttt{set} or \texttt{map}, elements are stored in a tree, not an array. \textbf{You cannot do \texttt{it + 5}}. You must step through them using \texttt{it++} or \texttt{++it}.

\subsection{Iterators and Pairs (\texttt{->} Operator)}
When an iterator points to a \texttt{pair} (which is what a \texttt{map} stores), dereferencing looks slightly different. You use the arrow operator (\texttt{->}):

\begin{lstlisting}
map<string, int> m;
m["Alice"] = 10;

auto it = m.begin(); // Iterator points to a pair<string, int>

// These two lines do the exact same thing:
cout << (*it).first << " " << (*it).second << "\n";
cout << it->first << " " << it->second << "\n"; // Much cleaner!
\end{lstlisting}

\vspace{0.5cm}
\hrule
\vspace{0.5cm}

\section{Homework}
Target the \textbf{CSES Problem Set}:
\begin{enumerate}
    \item \textbf{Introductory Problems:} Complete the entire section to solidify basic implementation.
    \item \textbf{Sorting and Searching (Selected based on today's lecture):}
    \begin{itemize}
        \item \textit{Distinct Numbers} (Apply \texttt{set} or \texttt{sort})
        \item \textit{Sum of Two Values} (Apply \texttt{map} or sorting + \texttt{lower\_bound})
        \item \textit{Restaurant Customers} (Apply \texttt{vector<pair<int, int>>} and custom sorting)
        \item \textit{Nearest Smaller Values} (Apply \texttt{stack})
    \end{itemize}
\end{enumerate}

\end{document}